\documentclass{situationist-preprint}

\title{The Stoic and the Problem of Priam}
\msid{ETH-11-b}
\msfamily{ETH-11}
\msdiscipline{Normative and Applied Ethics}
\mstype{Consequence}
\mskeywords{Stoicism; Priam; sufficiency of virtue; Antiochus; wretchedness; eudaimonia}

\begin{document}
\maketitle

\begin{abstract}
ETH-11-a argues that virtue is necessary but not sufficient for happiness, and that the case of Priam refutes the Stoic thesis. We draw out what this leaves the Stoic. The Stoic can keep the sufficiency thesis only by denying that Priam was virtuous, by denying that his losses were bad for him, or by identifying happiness with virtue; we show that each option either changes the subject or abandons the claim that happiness is a life that goes well. What survives is a weaker thesis, which Aristotle himself accepted: virtue is sufficient for a life that is not wretched. We show that the conclusion of ETH-11-a, with one further premise, yields this thesis, and that the compromise of Antiochus, on which virtue suffices for a happy but not the happiest life, collapses into it.
\end{abstract}

\section{The Stoic's options}

ETH-11-a defines happiness as a life that, taken as a whole, goes well to a high degree for the one who lives it (Definition~2.1), and virtue as possession of the virtues of character with practical wisdom to the degree attainable by a good human being (Definition~2.2). Its argument against sufficiency (Theorem~3.3) uses the principle that the goods of relationships, achievement and contact with reality are goods of a life (Principle~P1), the principle that what one has reason to prefer for one's own sake is good for one (Principle~P2), and the case of Priam (Assumption~3.2). The Stoic who wishes to retain the sufficiency thesis must reject one of these. There are three ways of doing so.

\paragraph{Denying that Priam was virtuous.} The Stoics held that virtue is a perfect and unified condition that few or none attain, and that all who lack it are equally vicious \citep{LongSedley1987}. On this conception Priam, a king of conspicuous piety and justice, was not virtuous, and his misery is no counterexample. But the conception gives the sufficiency thesis no human instances. Definition~2.2 of ETH-11-a concerns virtue to the degree attainable by a good human being, and the question whether virtue is sufficient for happiness, so understood, is the question that bears on how human beings should live. If the Stoic answers only that a perfect sage, should one exist, would be happy, he has not answered it. Moreover, the case can be rebuilt for any agent of attainable virtue, since nothing in Assumption~3.2 of ETH-11-a depends on Priam rather than another good person who loses what she loves.

\paragraph{Denying that the losses were bad for him.} This is the doctrine of preferred indifferents, discussed as Objection~1 of ETH-11-a. Its difficulty is internal. The Stoic grants that Priam had reason to select the lives of his sons for their own sake as well as for the sake of virtuous activity. By Principle~P2 of ETH-11-a, that reason makes their lives good for him. The Stoic must therefore deny Principle~P2, and with it the link between rational preference and good, which is the link on which the Stoic account of selection itself depends \citep{Annas1993}.

\paragraph{Identifying happiness with virtue.} The Stoic can say that \emph{eudaimonia} simply is the virtuous life, so that the virtuous person is happy by definition. Then the sufficiency thesis is true and empty: it no longer says anything about how well a life goes for the one who lives it. Definition~2.1 of ETH-11-a is not a stipulation that the Stoic may decline. It states what the ancient debate was about, as the Stoic's own appeal to happiness as the goal of life shows: the goal was meant to be what makes a life worth choosing for the one who lives it \citep[V]{CiceroTusc}.

\section{What survives}

Aristotle, having granted that one who meets with Priam's fortunes cannot be called blessed, adds that the virtuous person will never become wretched, since he will never do hateful and mean acts \citep[I.10, 1100b33--1101a8]{AristotleNE}. The thesis is weaker than sufficiency, and it is compatible with the conclusion of ETH-11-a.

\begin{assumption}[Floor]\label{a:floor}
Virtuous activity is a component of well-being, and a life in which it is sustained throughout does not, whatever else befalls it, go badly to a high degree for the one who lives it.
\end{assumption}

\begin{proposition}\label{p:floor}
Given Assumption~\ref{a:floor}, virtue is sufficient for a life that is not wretched, though, by Theorem~3.3 of ETH-11-a, not for a happy one.
\end{proposition}

\begin{proof}
A virtuous agent exercises the virtues in whatever circumstances she meets, including loss: Priam's courage and dignity in supplicating Achilles for his son's body are exercises of virtue in extremity. By Assumption~\ref{a:floor}, her life does not go badly to a high degree, so it is not wretched. Theorem~3.3 of ETH-11-a shows that it may still fail to be happy.
\end{proof}

Antiochus of Ascalon held that virtue is sufficient for a happy life but not for the happiest one, to which bodily and external goods contribute \citep[V]{CiceroFin}. The compromise faces a dilemma. If the happy life it promises meets the threshold of Definition~2.1 of ETH-11-a, Priam refutes it as he refutes the Stoic. If it does not, the happy life of Antiochus is a life above wretchedness and below happiness, and the compromise is Proposition~\ref{p:floor} under another name.

\section{Objection and reply}

\begin{objection}
The Stoic thesis is a practical doctrine. Its point is to direct the agent's concern to what is in her control. Refuting it as a theory of well-being leaves its point untouched.
\end{objection}

\begin{reply}
We agree that the practical doctrine survives, and ETH-11-a agrees that virtue is what is in the agent's control. But the doctrine is sound only if it is not supported by a falsehood. Telling Priam that he has lost nothing good is not consolation but misdescription; telling him that he has lost much, that his virtue remains his own, and that his life is not wretched, is true, and Proposition~\ref{p:floor} shows how it can be.
\end{reply}

\section{Conclusion}

The case of Priam leaves the Stoic three ways to keep the sufficiency thesis, and each either changes the subject or gives up the link between rational preference and good. What the Stoic can keep is the Aristotelian floor: virtue suffices for a life that is not wretched. That is the true thesis in the neighborhood of the Stoic one, and it is consistent with the conclusion that virtue is necessary but not sufficient for happiness.

\begin{thebibliography}{9}\small
\bibitem[Annas(1993)]{Annas1993} Annas, J. (1993). \emph{The Morality of Happiness}. Oxford University Press.
\bibitem[Aristotle(NE)]{AristotleNE} Aristotle. \emph{Nicomachean Ethics}. Cited by book, chapter and Bekker page.
\bibitem[Cicero(Fin.)]{CiceroFin} Cicero. \emph{De Finibus Bonorum et Malorum}. Cited by book.
\bibitem[Cicero(Tusc.)]{CiceroTusc} Cicero. \emph{Tusculan Disputations}. Cited by book.
\bibitem[Long and Sedley(1987)]{LongSedley1987} Long, A.~A. and Sedley, D.~N. (1987). \emph{The Hellenistic Philosophers}. Cambridge University Press.
\end{thebibliography}

\end{document}
